\subsection{A finite norm code does not recover a typed source}
\label{sec:grammar-fibers-12}

The refinement tests a concrete possible shortcut through the source
grammar. On the formal alphabet $\{1,2,3,4,6,8,9,12\}$, the unordered
triples with repetition number $\binom{10}{3}=120$. Their integer squared
norms have 107 classes, with thirteen two-member fibers. For example
$(1,6,8)$ and $(2,4,9)$ both have norm square 101 and sum 15, but products
48 and 72. Hence even sum plus squared norm is not an injective code.
Norm plus product gives 120 distinct signatures on this particular finite
dictionary.\samref{REF12_FOUR}

The complete equal-norm fibers are
\begin{center}\small
\begin{tabular}{rll}\toprule
Norm square&First triple&Second triple\\\midrule
41&(1,2,6)&(3,4,4)\\81&(1,4,8)&(3,6,6)\\
89&(2,2,9)&(3,4,8)\\101&(1,6,8)&(2,4,9)\\
146&(1,1,12)&(1,8,9)\\149&(1,2,12)&(2,8,9)\\
154&(1,3,12)&(3,8,9)\\161&(1,4,12)&(4,8,9)\\
164&(2,4,12)&(6,8,8)\\181&(1,6,12)&(6,8,9)\\
209&(1,8,12)&(8,8,9)\\226&(1,9,12)&(8,9,9)\\
289&(1,12,12)&(8,9,12)\\\bottomrule
\end{tabular}
\end{center}
These formal 120 candidates are not identified with the differently
selected 120 tested three-body rows earlier in the chapter merely by
cardinality. Neither code determines source type, admission, stability
or label reveal, and no universal triple-reconstruction theorem follows.
The witness instead fixes what must accompany a serialized source: its
typed construction and retained role/phase information.
